\documentclass[aspectratio=169,10pt,t]{beamer}
\usepackage{ac1}
\aulainfo{30}{Especificação do Projeto}{Do problema à estrutura do programa}

\begin{document}
\capa

\begin{frame}{Objetivo da aula}
  Cada grupo entrega uma \textbf{especificação} do projeto: descrição detalhada do problema, das funcionalidades principais e da estrutura prevista do programa.
  \vspace{1ex}
  \begin{itemize}
    \item Feedback \textbf{individualizado} para cada grupo
    \item Avaliação da \textbf{viabilidade}: o escopo cabe nas aulas restantes?
    \item Ajustes para que o projeto fique adequado ao nível da turma
  \end{itemize}
  \begin{exampleblock}{Lembrete do Cap.~9.1}
    ``Pense antes de programar'': uma boa especificação economiza horas de retrabalho.
  \end{exampleblock}
\end{frame}

\begin{frame}{Roteiro do documento de especificação (1 a 2 páginas)}
  \begin{enumerate}
    \item \textbf{Título e integrantes}
    \item \textbf{Problema}: o que será resolvido, para quem, por quê
    \item \textbf{Dados}: entradas (usuário, arquivos), saídas (tela, arquivos, gráficos)
    \item \textbf{Funcionalidades}: lista numerada, em ordem de prioridade
    \item \textbf{Estrutura do programa}: módulos e funções previstos
    \item \textbf{Estruturas de dados}: listas, dicionários, tuplas, arrays, DataFrames
    \item \textbf{Casos de teste e casos-limite}: o que pode dar errado?
    \item \textbf{Divisão de tarefas e cronograma}
  \end{enumerate}
\end{frame}

\begin{frame}{Exemplo: Gestão de insumos e custos (Eng.\ Civil/Produção)}
  \small
  \begin{columns}[T]
    \column{0.52\textwidth}
    \textbf{Problema.} Pequena obra precisa controlar o estoque e o custo de insumos.\\[1ex]
    \textbf{Entradas.} \texttt{insumos.csv} (nome, unidade, preço, estoque); comandos do usuário.\\[1ex]
    \textbf{Saídas.} Relatório de custos, alerta de estoque baixo, gráfico dos 5 insumos mais caros.
    \column{0.45\textwidth}
    \textbf{Funcionalidades}
    \begin{enumerate}
      \item Cadastrar/editar insumo
      \item Registrar entrada e saída de estoque
      \item Calcular custo total e por categoria
      \item Alertar quando estoque \(\le\) mínimo
      \item Salvar e carregar do arquivo CSV
      \item Gerar gráfico (Matplotlib)
    \end{enumerate}
  \end{columns}
\end{frame}

\begin{frame}[fragile]{Exemplo: estrutura prevista}
\begin{lstlisting}[basicstyle=\ttfamily\scriptsize]
projeto_insumos/
    principal.py        # menu e fluxo geral (main)
    estoque.py          # regras: cadastra(), baixa(), custo_total(), alertas()
    arquivos.py         # le_csv(), salva_csv()
    graficos.py         # grafico_mais_caros(), grafico_por_categoria()
    insumos.csv         # dados de exemplo
    testes.py           # asserts das funções de estoque
\end{lstlisting}
  \textbf{Estrutura de dados:}
\begin{lstlisting}[basicstyle=\ttfamily\scriptsize]
estoque = {
    "cimento": {"unidade": "saco", "preco": 38.5, "qtd": 120, "minimo": 30},
    "areia":   {"unidade": "m3",   "preco": 95.0, "qtd": 15,  "minimo": 10},
}
\end{lstlisting}
  \small Dicionário de dicionários (Aula 20): chave = nome do insumo; busca direta e legível.
\end{frame}

\begin{frame}{Exemplo para Eng.\ Elétrica/Mecânica}
  \small
  \begin{block}{Análise de ensaio de um resistor}
    \textbf{Entradas:} arquivo CSV com tensão e corrente.\\
    \textbf{Funcionalidades:} (1) ler e validar os dados; (2) ajustar a reta \(I=aV+b\) com NumPy; (3) estimar a resistência \(R=1/a\) e a incerteza; (4) detectar medidas discrepantes; (5) gráfico dos pontos com a reta; (6) relatório em arquivo texto.\\
    \textbf{Casos-limite:} menos de 2 pontos, tensão repetida, arquivo vazio, valores não numéricos.
  \end{block}
  \begin{block}{Séries históricas (Eng.\ Ambiental/Cartográfica)}
    \textbf{Entradas:} CSV de temperatura diária.\\
    \textbf{Funcionalidades:} médias mensais e anuais com Pandas; anos mais quentes; gráfico de tendência; conversor de coordenadas.
  \end{block}
\end{frame}

\begin{frame}{Como especificar boas funcionalidades}
  \begin{columns}[T]
    \column{0.48\textwidth}
    \begin{alertblock}{Vago (evite)}
      \begin{itemize}\small
        \item ``O programa gerencia tudo''
        \item ``Faz uns gráficos''
        \item ``Trata os erros''
      \end{itemize}
    \end{alertblock}
    \column{0.48\textwidth}
    \begin{exampleblock}{Verificável (prefira)}
      \begin{itemize}\small
        \item ``Mostra o custo total e o de cada categoria, com 2 casas decimais''
        \item ``Gera um gráfico de barras com os 5 itens mais caros e salva em PNG''
        \item ``Se o arquivo não existe, avisa e cria um novo''
      \end{itemize}
    \end{exampleblock}
  \end{columns}
  \vspace{1ex}
  \textbf{Teste:} para cada funcionalidade, escreva um caso de uso com entrada e resultado esperado.
\end{frame}

\begin{frame}{Planejando o tamanho do projeto}
  \begin{itemize}
    \item Classifique as funcionalidades: \textbf{essenciais} (sem elas não há programa), \textbf{desejáveis} e \textbf{extras}
    \item Restam \textbf{5 aulas} (31 a 35): implemente primeiro o essencial
    \item Estime: ``quantas funções? quantas linhas?'' (um projeto típico: 150 a 400 linhas)
    \item Evite \emph{interfaces gráficas} e \emph{bancos de dados}: fora do escopo da disciplina
  \end{itemize}
  \begin{alertblock}{O ótimo é inimigo do bom}
    Prefira um programa pequeno, correto e bem organizado a um grande que não funciona.
  \end{alertblock}
\end{frame}

\begin{frame}{Atividade desta aula}
  \begin{enumerate}
    \item Cada grupo escreve (ou finaliza) o documento de especificação seguindo o roteiro
    \item O professor atende os grupos um a um e verifica:
    \begin{itemize}\small
      \item O problema está claro? Dá para testar?
      \item O escopo é viável? O que é essencial?
      \item A divisão em funções e dados faz sentido?
    \end{itemize}
    \item Ajustem a especificação com base no feedback
    \item \textbf{Entrega ao fim da aula} (ou até a próxima, conforme combinado)
  \end{enumerate}
\end{frame}

\begin{frame}{Checklist da especificação}
  \small
  \begin{tabular}{@{}cl@{}}
    $\square$ & O problema cabe em duas frases?\\
    $\square$ & Entradas e saídas estão definidas, com formato dos arquivos?\\
    $\square$ & Há pelo menos 5 funcionalidades verificáveis, priorizadas?\\
    $\square$ & Há uma lista de funções/módulos e a estrutura de dados escolhida?\\
    $\square$ & Os casos-limite e as entradas inválidas estão listados?\\
    $\square$ & Cada integrante tem tarefas definidas?\\
    $\square$ & Os dados de exemplo já estão disponíveis?\\
  \end{tabular}
\end{frame}

\end{document}
